\chapter{Probability 101}
\label{chap:prob}

\vspace*{2cm}

\begin{center}
	\itshape
	"It is unanimously agreed that statistics
	depends somehow on probability. But, as to
	what probability is, and how it is connected
	with statistics, there has seldom been such
	complete disagreement and breakdown of
	communication since the Tower of Babel.
	
	
	Doubtless, much of the disagreement is
	merely terminological and would disappear
	under sufficiently sharp analysis."
	
	\vspace{0.5em}
	--- Savage, 1954, The foundations of statistics \cite{savage1972introduction}
\end{center}

\vspace{2cm}


Throughout this course we will proceed with a physical interpretation of Probability, as we approach this complex field through the eyes of experimental physicists who want to \emph{measure something}. In nature, there exists some measurable quantities whose value depends (in some way) from unknown parameters. Observing these observables $x\in X$ we can make inferences on the nature of these unknown parameters $\theta\in A$. 

In this first treatment we have introduced the notation that will follow us for the rest of the course - observable quantities will be denoted with the Latin letters \emph{x, y, z}, while unknown parameters will be denoted with the Greek letters \emph{$\theta, \mu, \nu, \sigma$}. The important subtlety that this hopes to highlight is that these quantities live in different spaces, and that observables are quantities that we can measure while the parameters are not directly measurable, and the best we can hope to do is make an inference based on a measurement. Let's never forget this, as it will be a fundamental to not mix the two quantities in the future. 



\section{Operative definition of Probability}
\label{sec:prob_def}

We shall start with an operative definition of Probability, consistent with a physical quantity. Given an experimental procedure that is infinitely replicable, each repetition will correspond to an observable quantity $x\in X$. We shall call \emph{event} $E$ the subset of possible outcomes of $x$, $E\subset X$. The \textbf{probability of event E} in the frequentist sense is defined as:
	\[
	P(E) = \lim_{N\to\infty} \frac{n(E)}{N}
	\]

where $n(E)$ is the number of times that event $E$ is verified out of $N$ trials. In this operative definition, $P(E)$ is a measurable quantity, and thus it is fundamental that the experiment be repeatable infinite times in this formalism.

$X$ is the space of the possible outcomes of the observable $x$, and in the simple discrete example of a mix of colored balls inside of a box $X=\{ \text{red, blue, green}\}$ and $E$ can be any single extraction of a ball, or any combination. Each event defined this way will have its own probability, and as long as the experiment is infinitely repeatable this probability well defined and given by the frequentist (Von Mises) definition provided above. 


However, as mathematical limits are, for a limit to be well defined it must equal the same value in whichever way $N$ decides to reach infinity, no matter how the events are selected. This means that if I repeat an experiment to infinity, the probability of event $E$ must be the same if I consider all of the events, just the even ones, just the prime numbered ones, etc. 

\begin{example}
	If in the previous situation I extract (with replacement), the colored balls randomly to infinity, the sequence 
	\[
	\text{red, blue, green, red, blue, green, red, blue, green, red, }\dots
	\]
	
	does not have $P(\text{red})$ well defined, as if I consider all of the events in the limit I get $P(\text{red})=\frac{1}{3}$, but if I consider all of the events that follow the extraction of a green ball, $P(\text{red})=1$ and all those following a blue extraction $P(\text{red})=0$. These three limits were all taken in the mathematical sense to infinity, though they are not equal to the same value and therefore the limit is not well defined, and in fact this deterministic sequence does not have a probability defined. 
\end{example}

This means that frequentist probability is defined on an ensemble.\todo{This should be looked upon further, or omitted}


\subsection{Kolmogorov probability}
\label{subsec:kolmogorov}

	We define "probability" on the set $S$ some function of the set $P$\sidenote{Formally, $P$ is defined on a $\sigma$-algebra of subsets of $S$, so that we are sure that the union and conjugation operations remain within the same space} with the following properties:
\begin{enumerate}
	
	\item $\forall A \subset S \quad P(A)\geq 0$
	\item $P(S) = 1$
	\item $A_i \subset S ;\quad A_i\cap A_j = 0 \quad \forall i\neq j ;\quad \forall A_i \rightarrow P(\cup A_i) = \sum\limits_{i}P(A_i)$
	
\end{enumerate}

$S$ can be taken as the space of all possible experimental observations, and each element a completely defined observation. In this case, $A$ represents a set of possible observations.


From these three axioms we can define a series of simple theorems\todo{Do we want to prove these? Also maybe the formatting could do with some work}:


\begin{definition}[Conditional probability]
	Probability of event A given that event B has occurred
	\[
	P(A|B) \equiv \frac{P(A\cap B)}{P(B)}
	\]
This can be re-written in the more useful form:
	\[
	P(A\cap B)= P(A|B)\cdot P(B)
	\]
\end{definition}

$P(A|B)$ (also denoted with $P_B(A)$) is a Kolmogorov probability and satisfies all of the necessary axioms\todo{Proof?}.

It is important to note that conditional probability has nothing to do with causation, correlation, or temporal relations between two events, simply put it is the probability of an event given another, and therefore is nothing more than a restriction of the probability space to a smaller ensemble. 

\begin{theorem}[Total probability law]
	Let $\{A_i\}$ be a countable partition of $S$ such that $\cup_i A_i = S$. It follows that:
	\[
	\forall B: P(B)=\sum_i P(B\cap A_i) = \sum_i P(B|A_i)P(A_i)
	\]
\label{th:total_prob_law}
\end{theorem}\todo{Here I omitted the intermediate passages for the proof which can be added in-line or as a proof later}

This theorem is useful when the conditional probabilities are easier to calculate than the total.
\begin{example}
	I know that $25\%$ of lightbulbs are produced in Vietnam, while $75\%$ are produced in China. I know that of the Vietnamese lightbulbs, one in a thousand are broken during production while the Chinese ones are broken one percent of the time. I buy a lightbulb from the store. What is the probability that it is broken?
	
	The total probability law comes to help, as we can write:
	\begin{align*}
	P(\text{Broken}) &= P(\text{Broken}|\text{C})\cdot P(\text{C}) + P(\text{Broken}|\text{V})\cdot P(\text{V}) \\
	&= 0.01\times0.75 + 0.001\times0.25 \\
	&= 0.775\%
	\end{align*}
\end{example}

\begin{definition}[Independent events]
	Two events are independent if 
	\[
	P(A|B)=P(A)
	\]
	or, reformulated, two events are independent if
	\[
	P(A\cap B) = P(A)\cdot P(B)
	\]	
	This second form takes the name of the product law of probability
\end{definition}

It is important to remember that independence is not correlation, and the former only describes whether or not the probability of an event changes after the verification of another. For example, rolling a die and flipping a coin are two independent events, while picking two cards from a deck without replacement are not two independent events as removing one card from the deck changes the odds of picking the second. 

\begin{theorem}[Bayes' theorem]
	We have already written that 
	\[
	P(B\cap A) = P(B|A)\cdot P(A) \overset{\scriptscriptstyle \text{symm.}}{=} P(A|B)\cdot P(B)
	\]
	It follows naturally that
	\[
	P(B|A) = \frac{P(A|B)\cdot P(B)}{P(A)}
	\]
\label{th:Bayes}
\end{theorem}	

The last expression takes the name of \textbf{Bayes' theorem}. Interestingly enough, it is not what characterizes Bayesian statistics, and we remind that we are still in the Kolmogorov frequentist approach, though the theorem generalizes to the Bayesian approach if the probabilities are replaced with degrees of belief, which we will briefly introduce in the next section.

\begin{example}
	Before moving on, let's illustrate the concepts we have encountered so far by solving the famous \textbf{Monty Hall problem} (you have \emph{probably} heard of it at least once in your life).
	
	Imagine you are suddenly transported from where you're reading this text to the studio of a famous game show. The host, Mr.~Monty, explains the rules of the game: there are three closed doors in front of you. Behind one there is a brand new sports car, behind the other two there are goats. You choose one door; if you pick the right one, you win the prize behind it. Of course, Mr.~Monty knows where the car is hidden.
	
	From what we discussed before, you know that each door has the same probability $\frac{1}{3}$ of hiding the car. So you pick, say, door number one (it does not really matter which door you choose for what follows).
	
	Now Monty wants to spice things up for the audience. After your choice, he opens door number three, revealing a goat (to be clear: he cannot open the door you chose, nor the one hiding the car, revealing the prize immediately would not make for a great show, would it?).
	
	At this point, he asks you to decide: do you want to stick with door number one, or switch to door number two?
	At first glance, this may seem like a pointless mind game: you might think you still have the same one-third chance of winning. But reality, as often happens in probability, is a bit more subtle and can even feel paradoxical. So you take a moment, pull out your sketchbook in front of the crowd, and do some math.
	
	We can tackle the problem introducing the following different events:
	\begin{itemize}
		\item $C_i$: the car is behind door $i$
		\item $A_i$: Mr.~Monty opens door $i$, revealing a goat
	\end{itemize}
	with $i \in \{1,2,3\}$.
	
	To maximize your chances of winning, you want to compare
	\[
	P(C_1|A_3)\lesseqqgtr P(C_2|A_3),
	\]
	that is: given that Monty opened door three, is it better to stick with door one or switch to door two?
	
	Using Bayes' theorem (Theorem~\ref{th:Bayes}), we can compute
	\[
	P(C_1|A_3)=\frac{P(A_3|C_1)\,P(C_1)}{P(A_3)}
	\]
	As expected, the prior probabilities are
	\[
	P(C_i)=\frac{1}{3}
	\]
	
	The term $P(A_3|C_1)$ is the probability that Monty opens door three given that the car is behind door one. Since in this case both doors two and three hide goats, and Monty can choose either, we have
	\[
	P(A_3|C_1)=\frac{1}{2}
	\]
	Next, we compute the total probability of the event $A_3$ using the law of total probability (Theorem~\ref{th:total_prob_law}):
	\begin{align*}
		P(A_3)
		&=\sum_i P(A_3|C_i)P(C_i)\\
		&=P(A_3|C_1)P(C_1)+P(A_3|C_2)P(C_2)+P(A_3|C_3)P(C_3)
	\end{align*}
	
	We already know the first term. The second term corresponds to the case in which the car is behind door two: Monty cannot open door one (your choice) nor door two (it hides the car), so he must open door three. Hence
	\[
	P(A_3|C_2)=1
	\]
	
	The third term corresponds to the case in which the car is behind door three. Since Monty cannot reveal the car, he will never open door three, so
	\[
	P(A_3|C_3)=0
	\]
	
	Substituting everything, we obtain
	\[
	P(A_3)=\frac{1}{3}\left(\frac{1}{2}+1+0\right)=\frac{1}{2}
	\]
	
	Therefore,
	\[
	P(C_1|A_3)=\frac{\frac{1}{2}\cdot\frac{1}{3}}{\frac{1}{2}}=\frac{1}{3}
	\]
	and consequently
	\[
	P(C_2|A_3)=1-P(C_1|A_3)=\frac{2}{3}
	\]
	
	This means that switching doors doubles your probability of winning the car!
	This result may seem surprising at first. The key point is that Mr.~Monty does not act randomly: he knows where the car is, and his choice is conditioned on that knowledge. By revealing a goat, he is giving you information, which updates your posterior probabilities.
	So go ahead: close your sketchbook, switch doors every time, and head back to where you were studying driving your brand new sports car (well... with double the probability, at least).
\end{example}


\subsubsection{Spoiler: Bayesian probability}
\label{subsubsec:Bayes}

The Bayesian school uses a concept of probability different than the frequentist (von Mises) concept. To try to minimize the natural confusion that will follow this chapter, we will baptize this concept with a different name: \emph{"degree of belief"} $\pi(x)$. The idea of this subjective reasoning is that events of which we have similar information are by definition equally probable. This concept, differently than the frequentist definition we have studied already, can be applied to whatever quantity, even if the latter is not observable. It does not depend on the ensemble, but on the subject itself, and Bayesian d.o.b (degree of belief) has been shown to satisfy the Kolmogorov axioms. We will be returning on this concept in various chapters\todo{Maybe add which ones through a ref, and perhaps elaborate further once we understand better} and it will be crucial to not confuse it with the frequentist definition of probability. 

\subsection{Events as physical observables}
\label{subsec:observables}

In physics, we usually call "event" a measurement of a certain observable $X$ yielding the value $x\in O_X$, where $O_X$ is the space of all the possible results (values of $x$). 


We classify $X$ as a \textbf{random observable} if the probability of measuring each of its values (or set of values\sidenote{Let's never forget that the probability of measuring an exact value from a continuous distribution is exactly zero, and therefore for continuous observables we will always talk about the probability of measuring a set of values, for example $0\leq x\leq 1$.}) exists and is well defined. 

Physically, this implies that for any subset $A \subset O_X$, the frequentist limit converges:

	\[
	\forall A \subset O_X \quad \exists P(x\in A) = \lim_{N\to\infty}\frac{n(x\in A)}{N}
	\]


To move beyond empirical observation, math comes to the rescue, and we map this physical intuition to the formal mathematical framework of measure theory. This allows us to treat "measurements" as "functions" and inherit a powerful toolkit of existing theorems.

\begin{definition}[Random variable]
	A random variable is a function $X: S\rightarrow O_X$ such that:
	\[
	\forall A \subset O_X \quad X^{-1}(A)\subset \Sigma \rightarrow \exists P(X^{-1}(A))
	\]
\end{definition}


The function $X$  maps the quantity we have observed to a value that we will call "random variable". Since the outcome of an experiment can be a qualitative observable, or anything really, this bridge is helpful in order to import all of the nice properties that mathematics guarantees us. However, for this function to be mathematically sound, we must confirm that it is measurable. That means that for any numerical interval we are interested in (say, the probability of a value falling between 5 and 10) there must exist a corresponding, well-defined set of real-world outcomes ($X^{-1}(A)$) that sits within our "event space" ($\Sigma$). If this condition is met, we gain the power to assign a precise probability $P$ to those outcomes.

In essence, the random variable ensures that every "observable state of the world" captured by our instruments has a mathematically guaranteed seat at the table of probability theory.

This may all sound scary right now, but hopefully the following example will clarify that the concept, at its core, is really quite simple. 

\begin{example}
	Consider the simplest physical experiment: flipping a coin. The Sample Space $S$ contains two physical outcomes: $\{\text{Heads,}$ $\text{Tails}\}$. To perform calculations, we define a Random Variable $X$ as a function that maps these physical states to numbers:
	
	\[
	X(s) = \begin{cases} 1 & \text{if } s = \text{Heads} \\ 0 & \text{if } s = \text{Tails} \end{cases}
	\]
	
	Here, the "Observable Space" $O_X$ is $\{0, 1\}$. If we want to find the probability of the set $A = \{1\}$ (the event "Heads"), we look at the pre-image:
	
	
	\[
	X^{-1}(\{1\}) = \{\text{Heads}\}
	\]
	
	Since "Heads" is a valid element of our event space $\Sigma$, the probability $P(X^{-1}(\{1\}))$ is well-defined (e.g., $0.5$). The function $X$ has successfully translated a physical piece of metal landing on a table into a number we can use in an equation.
\end{example}

Since we are looking to do \emph{something} with these random variables, it will surely be helpful to note that any function of a random variable is a random variable itself, and more importantly if a random variable assumes a probability distribution, each function of that variable will as well. Since we will be using "functions of random variables" quite often, let's give them a name: \textbf{statistics}.
\begin{definition}[Statistic]
	A statistic (singular) is a function of random variables $s(x,y,\dots)$, that depends only on the values themselves, not on unknown parameters. 
\end{definition}

The second part of this definition is important - a statistic only depends on \emph{measurable quantities} and is therefore also a measurable quantity. A nice, simple example of an important statistic is the mean of a set of data. This concept of statistic will follow us closely throughout the whole course, and while it's a good idea to introduce them now this will certainly not be the last time we update our definition.

\subsubsection{Probability distributions}
\label{subsubsec:pdf}

We have already thrown out this term \emph{probability distribution} in the previous section, and we won't get very far without giving it the proper definition it deserves. 

\begin{definition}[Probability mass function \emph{pmf}]
	For a discrete random variable X
	\[
	pmf(x)\equiv P(E:X=x)
	\]
\end{definition}

Often this gets (somewhat improperly) condensed in literature to be written as $P(X)$, which we will adopt from now on without too much hesitation. Probability mass functions are additive and satisfy our lovely Kolmogorov axioms\sidenote{They are defined positive and less than $1$, the sum of the total possibility space is $1$, and they are additive. The sum operation naturally changes definition based on the dimension of the outcome space.}, and so we can truly call them "probabilities". A simple example of a \emph{pmf} would be the case of the outcome of a fair six-sided die roll
	\[
	pmf(x) = \begin{cases} \frac{1}{6} & \text{if } x \in \{1,2,3,4,5,6\} \\ 0 &  \text{otherwise} \end{cases}
	\]


So much for discrete random variables, let's talk about the interesting stuff and generalize this formalism to continuous\sidenote{Naturally, we will never have a continuous result from an experiment as everything at some scale gets discretized or binned in some way, though we are sure the readers will have no problem with forgiving this abstraction.} observables. The probability mass function, as we have defined it, will be identically zero for each value in a continuous space such as $\Re$, so we will need a new definition, and much like we are used to in physics, the concept of \emph{mass} gets translated to \emph{density}.

\begin{definition}[Probability density function pdf]
	For a continuous random variable x
	\[
	P(x:x\in A)= \int_{A} p_X(x)\dd x
	\]
	where $p_X(x)$ is the \emph{probability density function} of $x$ 
\end{definition}\todo{If we find a better definition, let's be sure to put it here. This one works, but it's not satisfying}

\todo{Maybe we should include here a brief discussion of pdf's properties, e.g. how they add up, ecc.}

We use the term "function", but if we intend the integral in the Lebesgue sense we can extend the definition to \emph{pdf's} that are not functions but rather distributions. In general, this allows us a series of liberties that functions formally would not allow, such as mixing discrete and continuous variables. It's not guaranteed that a pdf satisfies the Kolmogorov axioms, and the following example shows that while $p(x)\geq 0$ not necessarily $p(x) < 1$.

	\[
	p(x) = \begin{cases} 2 & \text{if } x \in [0,\frac{1}{2}] \\ 0 &  \text{otherwise} \end{cases}
	\]


However, it is crucial that $\int p(x)\dd x = P(X\in O_X) = 1$ (we say that the pdf is \emph{normalized}), and one can verify that the example above satisfies this.\sidenote{Interesting to note as well that while probabilities in the Kolmogorov sense are pure dimensionless numbers, \emph{pdf's} have the dimensions of $[x]^{-N}$ where $N$ is the dimension of the space. This has natural parallels to the concept of mass and density in physics.} This normalization also brings us to define the \textbf{cumulative density function} of a distribution, as assures us that it is a monotonically increasing function of $x$ from $0$ to $1$.

\begin{definition}[Cumulative density function \emph{cdf}]
	For a continuous random variable x
	\[
	F_X(x)\equiv P\left(X\in \left(-\infty, x\right]\right) = \int_{-\infty}^{x}p_X(x^\prime)\dd x^\prime
	\]
\end{definition}

The fundamental theorem of calculus also tells us that $\dv{F}{x}=p(x)$. We will find the cumulative easier to calculate than the density function itself in many cases. 

An interesting and somewhat unexpected result can be reached when talking about \emph{CDF's}: the \textbf{probability integral transform} \cite{david1948probability} (also known as the \emph{universality of the uniform}\sidenote{If you haven't done so already, take a moment to brush up on what a Uniform distribution is, by heading over to the appendix \ref{sec:uniform}}).

\begin{theorem}[Probability integral transform]
	Suppose that a random variable $X$ has a continuous distribution for which the cumulative distribution function (CDF) is $F_X$. Then the random variable $Y$ defined as $Y \equiv F_X(X)$ follows a standard uniform distribution
	\[
	p_Y(Y) = U[0,1]
	\]
\label{th:universal_uniform}
\end{theorem}

\begin{example}
Let's prove this theorem:
Given a random continuous variable $X$, and $Y$ defined as above ($Y \equiv F_X(X)$) where $F_X$ is the \emph{cdf} of $X$. Given $y\in[0,1]$

\begin{align*}
	F_Y(y) &= P\left(Y\leq y\right) = P\left(F_X(X)\leq y\right)
\end{align*}

Since $F_X$ is continuous and strictly increasing on the support of $X$, its inverse $F_X^{-1}$ exists and is also increasing. Therefore,
\[
F_X(X) \le y
\quad \Longleftrightarrow \quad
X \le F_X^{-1}(y)
\]

Hence:
\begin{align*}
	F_Y(y) &= P\left(X\leq F_X^{-1}(y)\right)\\
	&= F_X\left(F_X^{-1}(y)\right) = y
\end{align*}

$F_Y(y) = y$ is the cumulative of a standard Uniform distribution, and so $Y$ is distributed as such.
\end{example}

This theorem will be useful in the far future when we talk about $p$-values, though it is interesting to explore the consequences that it suggests.

Moving on to probability distributions for more than one random variable, things generalize quite easily. We can generally call $P(x_1, \allowbreak x_2, \allowbreak \dots, \allowbreak x_N)$ the probability distribution of a countable number of random variables, each defined on its own outcome space. We can define independence as we have in section \ref{subsec:kolmogorov}, and say that probability distributions (defined on the same $S$) are independent if 


	\[
	\forall A,B \quad P\left( \left( x\in A\right) \cap \left( y\in B\right) \right)  = P\left( x\in A\right) \cdot P\left( y\in B\right) 
	\]
	


This can also be written as $P_{XY}(x,y)=P_X(x)\cdot P_Y(y)$.\todo{This condition is necessary, but not sufficient, so when you have an example, put it here}



Much like the concept of independence, we generalize conditional probability and Bayes' theorem to distributions, obtaining
	\[
	P_X(x|y)\cdot P_Y(y) = P_{XY}(x,y) = P_Y(y|x)\cdot P_x(x)
	\]
	
Finally, to end our parallels with Kolmogorov probability we consider a partition $\{Y_i\}$ of $O_Y$, particularly one where $Y$ is discrete-valued such that $Y_i = \{y_i\}$. In this case we have
	\[
	P_X(x) = \sum_i P_{XY}(x,y_i)=\sum_i P_X(x|y_i)\cdot P_Y(y_i)
	\]
	
where astute readers will have identified the total probability law. This process is called \emph{marginalizing}, and $P_X(x)$ is the \emph{marginalized} distribution of $P_{XY}(x,y)$ with respect to $y$. This concept, although easy, has been known to be a "surprise question" at the exams that students tend to overlook.\sidenote{One of the authors speaks from experience.}

\subsection{The law of large numbers}
\label{subsec:LLN}
The law of large numbers (LLN) is one of the most fundamental and well-known theorems of this field. Let's start with a few definitions:
\begin{definition}[Arithmetic mean]
	Given a random variable $X$, we define \textbf{arithmetic mean} (or \emph{sample mean on N}) the statistic:
	\[
	\bar{x}_N = \frac{1}{N}S_N = \frac{1}{N}\sum_i^N x_i
	\]
\end{definition}

\begin{definition}[Mean]
	If $p_X(x)$ is the pdf of $X$, we call the \emph{mean} of the distribution (or the expected value of $X$) the integral (Lebesgue):
	\[
	\langle x\rangle\equiv \E[x] \equiv \int_{O_X} x\cdot p_X(x)\dd x
	\]
\label{def:mean}
\end{definition}

\begin{theorem}[Law of large numbers (LLN)]
	Let $X$ be a real-valued observable, and $p(x)$ its pdf. Let's suppose that the mean of $X$ exists and is finite:
	\begin{equation}
	\mu = \E[x]=\int_{O_X} x\cdot p_X(x)\dd x
	\label{eq:LLN_mean}
	\end{equation}
	We then have that 
	\begin{equation}
	\forall \varepsilon > 0: \lim_{N\to\infty} P\left( \left| \bar{x}_N - \mu\right| > \varepsilon\right) =0
	\label{eq:LLN}
	\end{equation}
\end{theorem}


Within this theorem we have defined an interesting property "under the rug", which is convergence in probability, which is nothing more than the limit in (\ref{eq:LLN}).\todo{Maybe, instead of defining it this half-assed way we could add a definition explicitly?}

If we had to write the law of large numbers in words, it would read: "the sample mean of an observable converges in probability to its expected value whenever the latter exists and is finite". \todo{Here, the professor adds some considerations that I don't think I fully understand. Let's discuss if we want to include them.}\sidenote{For those of you who are not on their first read of these notes, or have some prior knowledge on point estimators, we can say that "the arithmetic mean is a consistent estimator of the mean $\E[x]$ of the distribution."}

\begin{example}[Proof of the LLN]
	The most general proof of the LLN is beyond the scope of this course and therefore this write-up, however we will look at a more limiting but still instructive example. Let's suppose that also the following integral exists (\emph{variance of} $X$):
	\[
	\sigma^2 \equiv \Var[X] \equiv \int_{O_X}\left( x-\E[x]\right)^2p_X(x)\dd x
	\] 
\begin{lemma}[Čebyšëv's theorem]
	Let $X$ be a real-valued observable for which mean and variance $\mu$ and $\sigma^2$ exist and are well defined. In this case, the following inequality:
	\begin{equation}
	P\left( \left| x-\mu \right|>k\sigma \right)\leq \frac{1}{k^2}
	\label{eq:chebyshev}
	\end{equation}
	exists for each $k\in\R^+$. While it's not the most useful upper bound it's useful in the proof of the LLN. 
\end{lemma}

\begin{example}[Proof of Čebyšëv's theorem]
		\begin{align*}
			P\left( \left| x-\mu \right|>k\sigma \right) &= \int_{\left| x-\mu \right|>k\sigma} p_X(x)\dd x \\
		&= \int_{O_X} \mathscr{F}\left( \left| x-\mu \right|>k\sigma \right)p_X(x)\dd x
	\end{align*}
	We have introduced $\mathscr{F}$ as the indicator function of the set $X$:
	\[
	\mathscr{F}\left( \left| x-\mu \right|>k\sigma \right) = \begin{cases} 1 & \text{if } \left| x-\mu \right|>k\sigma  \\ 0 &  \text{otherwise} \end{cases}
	\]
	\begin{figure}[H]
		\centering
		\input{images/chebyshev_theorem.pgf}
		\caption{Should we add a caption and a label to this image?}
	\end{figure}
	So substituting into the integrand this property:
	\[
	\mathscr{F}\left[ \frac{\left| x-\mu \right|}{k\sigma} > 1 \right] \leq \frac{\left| x-\mu \right|^2}{k^2\sigma^2} 
	\]
	we obtain 
	\[
	P\left( \left| x-\mu \right|>k\sigma \right) \leq \int \frac{1}{k^2\sigma^2}\cdot \underbrace{\left| x-\mu \right|^2 p_X(x)\dd x}_{\text{Variance}} = \frac{\Var[x]}{k^2\sigma^2} = \frac{1}{k^2}
	\]
	\qed
\end{example}

Now we can get to proving the LLN, armed with the Čebyšëv inequality. We notice that the pdf of the vectorial observable $\{x_i\}_N$ can be written as the product of the \emph{pdf's} of the single observables, assuming that the observations are all independent: $p(x_N)=\prod_i^N p_X(x_i)$. 

Therefore, we have that
\begin{align*}
	\Var{\bar{x}_N} &= \Var\left[ \frac{\sum_i x_i}{N} \right] = \frac{\Var\left[ \sum_i x_i \right]}{N^2} \\
	&= \frac{1}{N^2}\int\left[ \sum_i x_i - \E\left[ \sum_i x_i \right] \right]^2\prod_i p_X(x_i)\dd x^N \\ 
	&= \frac{1}{N^2}\int\left[ \sum_i x_i - \sum_i \mu \right]^2\prod_i p_X(x_i)\dd x^N & \text{(linearity of $\E$)}\\ 
	&= \frac{1}{N^2}\int\left[ \sum_i \left( x_i - \mu \right) \right]^2\prod_i p_X(x_i)\dd x^N \\
	&= \frac{1}{N^2}\int \sum_{i,j}\left[\left( x_i - \mu \right)\cdot\left( x_j - \mu \right)\right]\prod_i p_X(x_i)\dd x^N\\
	&= \frac{1}{N^2}\sum_{i}\int \left( x_i - \mu \right)^2p_X(x_i)\dd x^N & \text{(independent events)}\\
	&= \frac{\sigma^2}{N}
\end{align*}

Let's also note that
	\[
	\langle x \rangle = \int \bar{x}_N \prod_i p_X(x_i)\dd x^N = \mu
	\]
	
Applying the Čebyšëv inequality (\ref{eq:chebyshev}) to the statistic $\bar{x}_N$, with $k=\frac{\varepsilon}{\sqrt{\Var[{\bar{x}_N}]}}$ we get:
	\[
	P\left( \left| \bar{x}_N-\mu \right| > \varepsilon\right) \leq \frac{\Var[{\bar{x}_N}]}{\varepsilon^2} = \frac{\sigma^2}{N\varepsilon^2} 
	\]
Regardless of $\varepsilon$, the right hand side can be made arbitrarily small with a large enough $N$, and so it immediately follows that
	\[
	\forall \varepsilon \lim_{N\to\infty} P\left( \left| \bar{x}_N-\mu \right| > \varepsilon\right)= 0 
	\]
	\qed
\end{example}\todo{Phew... Look over all that to make sure nothing can be made clearer}
\todo{We used here a property of the Var we have not yet stated. While trivial if you look at the Var definition it might be useful to have a list of properties for quick reference}

If that wasn't satisfying enough the reader will be pleased to know that there is a stronger version of the Law of Large Numbers, appropriately called "Strong-LLN". 

\begin{theorem}[Strong-LLN]
	\[
	P\left( \lim_{N\to\infty} \left| \bar{x}_N - \mu \right| > \varepsilon \right) = 0
	\]
\end{theorem}

At first glance, this looks identical to what we have previously stated, but the subtle difference lies in the "probability of the limit" rather than the "limit of the probability". We can state that the limit of each succession coverges to $\mu$ "almost certainly" (which is stronger than convergence in probability and it actually implies the latter.) The SLLN is harder to prove and we won't even use it in the course of this text, but readers are welcome to take the opportunity to deepen their understanding on the matter.\sidenote{Or happily skip it and enjoy one less topic to study} The important topic to take away is that it exists, and that this "almost certain" convergence implies the existence of the expected value $\E[x]$. We shall not touch on the topic further.

\section{Operations on statistics}
\label{sec:stat-operations}

We have, in the previous sections, talked about operators such as the distribution, expected value, and variance of random variables. It is now time to generalize these definitions to statistics, which we remind are \emph{functions} of the random variables themselves. 

\subsection{Probability distribution of a statistic}
Let's begin with the most pressing one: how do we define the \emph{distribution} of a \emph{statistic}? Well, if I know the discrete pmf of $x$ then the pmf of the statistic $f(x)$ is easily written as $pmf\left(x(f)\right)$, no surprises there. 

Things get more interesting in the continuous case, where the interval in the variable space is does not necessarily have the same measure as its image in the statistic space. As we well know, we must multiply by the derivative of the statistic with respect to the random variable. To be more precise:
\begin{definition}[Distribution of a statistic]
	Given $p_x(x) (x\in\R)$
	\[
	p_f(f) = p_x\left(x(f)\right)\left|\dv{x}{f}\right| = p_x\left(x(f)\right)\left|\dv{f}{x}\right|^{-1}_{x=x(f)} 
	\]
\label{def:pdf_var_change}
\end{definition}

This simplifies considerably if we consider a monotonically increasing $f$, in which case the cumulative function of $f$ becomes
	\[
	F_f(f)= F_x\left(x(f)\right)
	\]
taken with a "$-1$" before if $f$ is monotonically decreasing. The reader can readily verify that deriving the above expression one finds the one written in definition \ref{def:pdf_var_change}

If we consider, for example, the transformation $y=\frac{x}{c}$, where $c$ is a constant, we get:
\[
	F_y(y)=F_x(yc)
\]
Conversely, if $y=xc$:
\[
	F_y(y)=F_x\left(\frac{y}{c}\right)
\]
These two situations appear frequently in practice, so it is useful to keep these relations in mind.

\begin{example}
	Suppose we start from a known distribution in $x$, and want to find the \emph{function} $f$ that brings us to a known distribution $p_f(f)$. The, inverting \ref{def:pdf_var_change} we arrive at:
	\[
	\left| \dv{x}{f} \right| = \frac{p_f(f)}{p_x(x(f))}
	\]
	
	To give an example, let's suppose $p_x(x) = 1; \, x\in[0,1]$ and $p_f(f) = ke^{-kf}; \, f\geq 0$ (uniform and exponential distributions, respectively). In this case:
	\begin{align*}
	\left| \dv{x}{f} \right| = p_f(f) \Rightarrow x(f) &= \pm \int_0^f p_f(f^\prime)\dd f^\prime + x(0) \\
	&= \pm\left( 1-e^{-kf}\right) + \genfrac{\{}{\}}{0pt}{1}{0}{1}
	\end{align*}
	I choose the monotonically decreasing one:
	\[
	x(f) = e^{-kf} \Rightarrow f = -\frac{\ln{x}}{k}
	\]
\end{example}

\begin{example}
	Consider an observable $x \in \mathbb{R}$ with probability density function (pdf) $p_x(x)$. Define the statistic $y = x^2$. We want to determine the pdf $p_y(y)$.

	Since the transformation is not one-to-one, we have $x(y) = \pm \sqrt{y}$. Therefore, two points in the original space map to the same value of $y$ (see Fig.~\ref{fig:variable_change}), and we must sum their contributions.
	
	\begin{figure}[H]
		\centering
		\input{images/variable_change.pgf}
		\caption{Illustration of the variable transformation $y = x^2$ when $x \sim \mathcal{N}(0, 1)$.}
		\label{fig:variable_change}
	\end{figure}
	
	Using Definition~\ref{def:pdf_var_change}, we obtain:
	\begin{align*}
		p_y(y)&=
		p_x(\sqrt{y})\left|\frac{\dd (\sqrt{y})}{\dd y}\right|
		+
		p_x({\scriptstyle -}\sqrt{y})\left|\frac{\dd ({\scriptstyle -}\sqrt{y})}{\dd y}\right|\\
		&=\frac{p_x(\sqrt{y})+p_x({\scriptstyle -}\sqrt{y})}{2\sqrt{y}} \qquad y \ge 0
	\end{align*}
	
	\paragraph{Uniform distribution}
	
	Assume now that $x$ follows a uniform distribution on $[-1,1]$:
	\[
	p_x(x)=
	\begin{cases}
		\frac{1}{2} & x \in [-1,1],\\
		0 & \text{otherwise}
	\end{cases}
	\]
	Then:
	\[
	p_y(y)=
	\begin{cases}
		\frac{1}{2\sqrt{y}} & 0 < y \le 1,\\
		0 & \text{otherwise}
	\end{cases}
	\]
	which the reader can easily verify is correctly normalized.
	
	\paragraph{Gaussian distribution}
	Now let $x$ follow a standard normal distribution (as in Fig.~\ref{fig:variable_change}):
	\[
	p_x(x)=\frac{1}{\sqrt{2\pi}} e^{-\frac{x^2}{2}}
	\]
	Then:
	\begin{align*}
		p_y(y)&=\frac{1}{\sqrt{2\pi}}\frac{e^{-\frac{(\sqrt{y})^2}{2}}+e^{-\frac{({\scriptstyle -}\sqrt{y})^2}{2}}}{2 \sqrt{y}}\\
		&=\frac{1}{\sqrt{2\pi}}\frac{e^{-\frac{y}{2}}}{\sqrt{y}} \qquad y \ge 0
	\end{align*}
	This is a $\chi^2$ distribution with one degree of freedom. We will encounter it frequently in the rest of these notes.

\end{example}

\subsection{Expected value of a statistic}
Another important operator is the expected value. We defined it for a random variable in Def \ref{def:mean}, and we now look to generalize it to statistics.
	\[
	\E[f(X)] = \E_X[f] = \int f(x)p_X(x)\dd x \,\,\, \text{or} \,\,\, \sum_i f(x_i)P_X(x_i) \,\, \text{if X is discrete}
	\]
	
$\E_X$ is a linear operator, and for the LLN $\E_X[x]=\mu$. We also take the opportunity to introduce some helpful properties:
	\[
	\Var[X] = \E_x\left[ (x-\mu)^2 \right] = \E_X[x^2] - \mu^2
	\]

This gives us a means of calculating the Variance in a much easier way, and easily shows how $\Var[x]\geq0$.

We can also apply the familiar Schwartz inequality:

	\[
	\E[XY]\leq\sqrt{\E[X^2]\cdot\E[Y^2]} \quad \text{(Schwartz inequality)}
	\]

This type of formalism brings us to say something quite cool: the expected value of any statistic of the observable $X$ is independent of the base observables on which the probability itself is defined, and instead only depends on the probability in $S$, the space of all the elemental events. To put it mathematically:

	\[
	\E_X[f] = \int f(x)p_X(x)\dd x = \int fp_F(f)\dd f = \E_F[f] = \mu_f
	\]

Therefore, we don't lose any generality if we omit the index and just write $\E[f]$.


We physicists are actually quite used to this type of trick, we are not saying anything more complicated than the fact that the expected value of a statistic does not change if you change the variable you calculate it in. However, since these concepts can be confusing the first time they are studied, the following example should hopefully clear the air. 

\begin{example}
	Let's suppose we have a gas whose particles have a velocity that follow a certain probability density function on an interval $p_v(v)$. Our goal is to find the average kinetic energy of the gas. 
	
	For simplicity, we'll assume that the velocity of the gas follows the following distribution (which we will later learn to be the very important \emph{uniform distribution})
	\[
	p(v) = \begin{cases} \frac{1}{v_{max}} & \text{if } v \in [0,v_{max}] \\ 0 &  \text{otherwise} \end{cases}
	\]
	Proceeding via direct calculation:
	\[
	\E_v[K(v)] = \int K(v)\cdot p_v(v) \dd v = \int_0^{v_{max}} \frac{1}{2}mv^2 \cdot \frac{1}{v_{max}} \dd v = \frac{v_{max}^2m}{6}
	\]
	
	Otherwise, we can choose to calculate the expected value of the probability density function of the kinetic energy itself, and transforming $p_v(v)$ into $p_K(K)$ using definition \ref{def:pdf_var_change}:
	
	\[
	p_K(k) = p_v(v(K))\cdot \left| \dv{v}{k} \right| = \frac{1}{v_{max}}\cdot\frac{1}{\sqrt{2Km}}\bigg\rvert_{K\in(0,\frac{1}{2}mv_{max}^2]}
	\]
	
	we can now calculate the expected value of this distribution as we know how to
	\[
	\E_K[K] = \int K\cdot p_K(K) \dd K = \int_0^{\frac{1}{2}mv_{max}^2} K \cdot \frac{1}{v_{max}}\cdot\frac{1}{\sqrt{2Km}} \dd K = \frac{v_{max}^2m}{6}
	\]
	\qed
\end{example}

What happens if we try to find the expected value of the multiplication of two independent variables? Well, direct calculation tells us that given $F=f(X), G=g(X)$
\begin{align*}
	\E[fg] &= \int f\cdot g \cdot p_{FG}(f,g)\dd f \dd g \\
	&= \int f\cdot g \cdot p_F(f)\cdot p_G(g)\dd f \dd g \\
	&= \int f p_F(f) \dd f \cdot\int g p_G(g)\dd g = \E[f]\E[g]	
\end{align*}
Generalizing this concept we arrive quite naturally to the mystic \textbf{covariance}.
\begin{definition}[Covariance]
	\[
	\operatorname{cov}[f,g] = \E\left[ (f-\E(f))(g-\E(g))\right] = \E[fg] - \E[f]\E[g]
	\]
	\[
	\operatorname{cov}[f,f] \equiv \Var[f]
	\]
\end{definition}\todo{Should we define here the Cov matrix and Corr?}
We get that the covariance is null for independent variables, but not vice-versa, as the following example demonstrates.

\begin{example}
Let's consider the case where $F=f(X)=X$ and $G=g(X)=X^2$, and let $X$ follow a uniform distribution centered at zero:
	\[
	p(x) = \begin{cases} \frac{1}{2} & \text{if } x \in [-1,1] \\ 0 &  \text{otherwise} \end{cases}
	\]
(which the reader can easily verify as correctly normalized).\\
The covariance between the two is calculated as:
	\[
	\operatorname{cov}[f,g]=\E\left[X^3\right]-\E\left[X\right]\E\left[X^2\right]
	\]
since the distribution is centered at zero, we have that $\E\left[X\right]=\mu=0$, thus leaving us with:
	\[
	\operatorname{cov}[f,g]=\E\left[X^3\right]
	\]
However, because $x^3$ is an odd function, its expected value over a symmetric interval is zero. This results in a null covariance, even though the two variables are clearly not independent.
\end{example}

\subsection{Moments of a statistic}
\label{subsec:moments}
We can further generalize the concept of mean and variance by introducing the \textbf{moments of a distribution}.
\begin{definition}[$n$\emph{-th} order moment of a distribution]
	\[
	\mu_n(x_0) = \E\left[ (x-x_0)^n\right]
	\]
\end{definition}

We have already encountered two of these, and tacking on the banal $0$-order case we have:\sidenote{Careful of the argument of these moments, they are not all evaluated for different $x_0$!}
	\[
	\mu_0(x_0) = 1 \quad \mu_1(0) = \mu \quad \mu_2(\mu_1) = \Var[x]
	\]

To translate between "central" ($x_0=0$, often denoted as $\mu_n^\prime$) moments and generic ones we can do so using the linearity of $\E$ and the Newton binomial expansion: 
	\[
	\mu_n(x_0) = \sum_{k=0}^n \binom{n}{k} \mu_k(0)(-x_0)^{n-k}
	\]

Two notes: 
\begin{enumerate}
	\item Even functions have all odd central moments equal to zero
	\item Only $\mu_0(x_0) = 1$ is guaranteed to exist finite
\end{enumerate}

While the first condition is intuitive, as even distributions must be symmetric about their central moment by definition, the second one is important to deepen, and so an example is provided.\todo{Finish this example, I'm sure he does the calculation in some lecture so we can take it from there}

\begin{example}
	Cauchy (Lorentz) distribution:
	\[
	p(x) = \frac{1}{\pi} \frac{1}{1+x^2}
	\]
	proof that $\mu_1$ does not exist and that $\mu_2$ diverges...I'm pretty sure that also all the odd moments DNE? Check it!
	
	What this means is that the LLN is not applicable to Cauchy-distributed data!
\end{example}\todo{Based on my notes: For a distribution symmetric around zero, all odd central moments are zero (if they exist), while the even central moments may exist or diverge depending on the tails of the distribution.}

The moments are quite a powerful tool, as if I know them all I can actually univocally determine the starting distribution.\todo{The following theorem can be proven, or at least expanded upon further}
\begin{theorem}
	If $p_1(x)$ and $p_2(x)$ have the same moments 
	
	and $\left(p_1(x)-p_2(x)\right)$ is expandable in a power series then 
	\[
	p_1(x) = p_2(x)
	\]
\end{theorem}

As we have seen, knowing the moments of a distribution is a powerful tool that has numerous implications, though actually calculating them can be non trivial. Luckily, we physicists love our good, old Taylor approximations and this is a good a time as any to put it into practice.

\todo{are these mu's or thetas? BOH}
\begin{align*}
	f(\vec{x}) = f(\vec{\mu}) + &\frac{\partial f}{\partial x_i}(x_i-\mu_i)+\frac{1}{2}\frac{\partial^2f}{\partial x_i \partial x_j}(x_i-\mu_i)(x_j-\mu_j)+o\left(\left(\vec{x}-\vec{\mu}\right)^3\right) \\
	\Rightarrow \E[f(\vec{x})] &\simeq \E[f(\vec{\mu})] + \frac{\partial f}{\partial x_i}\cancelto{0}{\E[x_i-\mu_i]} \quad + \frac{1}{2}\frac{\partial^2f}{\partial x_i \partial x_j}\E[(x_i-\mu_i)(x_j-\mu_j)] \\
	&\simeq f(\vec{\mu}) + \frac{1}{2}\frac{\partial^2f}{\partial x_i \partial x_j}\operatorname{cov}[x_i,x_j]
\end{align*}

We now consider the linear approximation (where $\E[f(\vec{x})]\approx f(\vec{\mu})$ and $f(\vec{x}) \approx f(\vec{\mu}) + \frac{\partial f}{\partial x_i}(x_i-\mu_i)$). In this regime, we have:

	\[
	f(\vec{x}) - \E[f(\vec{x})] \approx \sum_i \frac{\partial f}{\partial x_i}(x_i-\mu_i)
	\]
	
To compute the covariance of two functions $f_1$ and $f_2$, we consider:

	\begin{align*}
		\operatorname{cov}[f_1,f_2]&=\E \left[(f_1(\vec{x})-\E[f_1(\vec{x})])(f_2(\vec{x})-\E[f_2(\vec{x})])\right] \\
		&\approx \E \left[\left(\sum_i\frac{\partial f_1}{\partial x_i}(x_i-\mu_i)\right)\left(\sum_j\frac{\partial f_2}{\partial x_j}(x_j-\mu_j)\right)\right] \\
		&= \sum_i \sum_j \frac{\partial f_1}{\partial x_1} \frac{\partial f_2}{\partial x_j} \E \left[(x_i-\mu_i)(x_j-\mu_j)\right]
	\end{align*}

Which directly leads to:
	\[
	\operatorname{cov}[f_1,f_2] \approx \sum_{i,j}\frac{\partial f_1}{\partial x_i}\frac{\partial f_2}{\partial x_j}\operatorname{cov}[x_i,x_j]
	\]
Which for uncorrelated $x$ (diagonal covariance matrix) simplifies to:
	\[
	\sigma_f^2 = \Var[f] = \sum_i \left| \frac{\partial f}{\partial x_i} \right|^2\sigma_i^2 +\dots
	\]
	
We recognize this as our familiar \textbf{formula for the propagation of errors}. Interesting to note that this formula comes from a Taylor approximation.\todo{I tried to rewrite this part to make it clearer. I don't know if I accomplished the goal}
	
\begin{definition}[Moment generating function]
	We define \emph{moment generating function} for a distribution 
	\[
	M(t) = \E[e^{tx}]
	\]
	If M(t) exists, is finite and differentiable near the origin, then the distribution has infinite moments given from
	\[
	\mu^\prime_n = \frac{\partial^n M}{\partial t^n}\bigg\rvert_{t=0} \quad \mu_n(x_0) = \frac{\partial^n \left( Me^{-tx_0} \right)}{\partial t^n}\bigg\rvert_{t=0}
	\]
	Where
	\[
	\frac{\partial^n M}{\partial t^n} = \frac{\partial^n}{\partial t^n} \E\left[e^{tx}\right] =  \E\left[\frac{\partial^ne^{tx}}{\partial t^n}\right] = \E\left[x^ne^{tx}\right]\bigg\rvert_{t=0} = \E\left[ x^n \right] = \mu_n^\prime
	\]
	So the partial derivatives of this generating function are the moments themselves (about the central)
\end{definition}
	
$M(t)$ has the nice property that for $x=x_1+x_2$ (with $x_1,x_2$ \emph{independent}) $M(t)=M_1(t)M_2(t)$. We see this right away if we explicitly write 
	\[
	M(t) = \E\left[ e^{t(x_1+x_2)}\right] = \E\left[ e^{tx_1}\right]\E\left[ e^{tx_2}\right] = M_1(t)M_2(t)
	\]
	
This gives us a nice way to find the resulting pdf\sidenote{Or, at least its moments!} of the sum of two or more observables.
\todo{There are some more properties and examples discussed in the slides that can be added here as we see fit, like the example that brings to the fact that the variance sums with the square.}

On a similar vein to the moment generating function we can define the \textbf{characteristic function} of a probability distribution:

\begin{definition}[Characteristic function]
	In the cases where the moment generating function is not defined, we introduce the characteristic function of a probability distribution:
	\[
	\phi(t) = \E\left[ e^{itx} \right]
	\]
	with
	\[
	\mu_n^\prime = \frac{\partial^n \phi}{\partial (it)^n}\bigg\rvert_{t=0} = (-i)^n\frac{\partial^n \phi}{\partial t^n}\bigg\rvert_{t=0}
	\]
\end{definition}

In general, this function has analogous properties as $M$, though it exists for a wider class of distributions. Interesting to note that the characteristic function is essentially the Fourier transform of the probability distribution.\todo{Paul Levy theorem, I have never once used this in my life. Should we include?} \\
\textcolor{red}{To the best of my understanding, Levy's theorem it's useful because allows us to conclude that: if the characteristic function converge to a finite limit (as in the CLT), than the corresponding distribution (or better the CDF) also converge to a finite limit. We could simply talk about the theorem in a strictly qualitative way, saying something along the line of: "The theorem tells us that the limits in the $\phi$-space behave in the same way as in the $F$-space".}

\section{Central Limit Theorem}
Now seems like a great time to introduce what is probably the most important theorem we will talk about in this whole text. 

\begin{theorem}[Central limit theorem]\label{th:central_limit}
	Let $x$ be a real observable with a distribution $p(x)$ that admits the first two moments. Therefore, both the expected value and variance exists finite and are: 
	
	\[
	\E[x] = \mu \quad \Var[x] = \sigma^2
	\]
	
	Then, let $\bar{x}$ be the sample mean:
	\[
	\bar{x} = \frac{1}{N} \sum_{i=1}^{N} x_i
	\]
	
	We define the statistic $y = x - \mu$, take its sample mean and multiply by the square root of the sample size. Doing so, we obtain the new N-dependent statistic:
	
	\[
	y_N = \sqrt{N}\left(\bar{x}-\mu\right)
	\]
	
	In the limit of large $N$, this statistic tends to a distribution of universal form, independent of the specific shape of $p(x)$, only its mean and variance.
	
	\begin{align*}
		p_{y_N}(y_N) \xrightarrow[N\to\infty]{} &\frac{1}{\sqrt{2\pi\sigma^2}}\exp\left(-\frac{y_N^2}{2\sigma^2}\right) \\
		= &\N\left(y_N;0,\sigma^2\right)
		\label{eq:CLT}
		\tag{Central limit theorem}
	\end{align*}
\end{theorem}

\begin{corollary}
	Knowing the distribution $p_{y_N}(y_N)$ we can, using a variable change, find the more useful distribution of the sample mean, which is what we were looking for in the first place:
	
	\begin{equation}
		\bar{x} = \frac{y_N}{\sqrt{N}} + \mu \sim \N\left(\bar{x};\mu,\frac{\sigma^2}{N}\right)
	\end{equation}
	
	This important result shows how the variance of the mean scales with the inverse of $N$
\end{corollary}

\begin{example}[Proof of the central limit theorem]

	We begin by noting that we already have $\E[y_N]=0$ and $\Var[y_N]=\sigma^2$ independently of $N$, as the variance is additive and goes with the square of any multiplicative constant. Using the properties of the characteristic function:
	\[
	\phi_N(t)\equiv \phi_{y_N}(t) = \phi_y^N\left(\frac{t}{\sqrt{N}} \right)
	\]
	from which we write
	\begin{align*}
		\log\left(\phi_N(t)\right) =& N\log\left[ \phi_y \left(\frac{t}{\sqrt{N}}\right) \right] \\
		\simeq& N\log\left[ 1 + \frac{\partial\phi_y}{\partial t}\frac{t}{\sqrt{N}} + \frac{\partial^2\phi_y}{\partial t^2}\frac{t^2}{2N} + o\left(\frac{t}{\sqrt{N}}\right)^2\right]\\
		=&N\log\left[1+0-\frac{\sigma^2t^2}{2N}+ o\left(\frac{t}{\sqrt{N}}\right)^2\right]\\
		\simeq& N \left(-\frac{\sigma^2t^2}{2N}+ o\left(\frac{t}{\sqrt{N}}\right)^2\right)\\
		=&-\frac{\sigma^2t^2}{2}+ N\cdot o\left(\frac{t^2}{N}\right)\\
		\Rightarrow\log\left(\phi_N(t)\right)\rightarrow&\log\phi(t)=-\frac{\sigma^2t^2}{2}
	\end{align*}
	
	Therefore:
	
	\[
	\phi(t) = \lim_{N\to\infty}\phi_N(t)=\exp\left(-\frac{\sigma^2t^2}{2}\right)
	\]
	
	This $\phi(t)$ is universal, and does not depend on $p(x)$
	
	We now want to recover the shape of this particular distribution whose characteristic function is $\phi(t)$. In other words, we want to perform a reverse Fourier transform to recover the explicit form of the limit distribution.
 
	\begin{align*}
		g(y_N)&=\frac{1}{2 \pi}\int_{-\infty}^{\infty}\exp\left(\frac{-\sigma^2t^2}{2}dt\right)\exp\left(-ity_N\right) dt\\
		&=\frac{1}{2 \pi}\int_{-\infty}^{\infty}\exp\left(\frac{-\sigma^2t^2-2ity_N-i^2y_N^2/ \sigma^2+i^2y_N^2/ \sigma^2}{2}\right) dt\\
		&=\frac{1}{2 \pi}\int_{-\infty}^{\infty}\exp\left(-\frac{\left(\sigma t -iy_N/\sigma \right)^2}{2}\right)\exp\left(-\frac{y_N^2}{2\sigma^2}\right) dt\\
		&=\frac{\exp\left(-\frac{y_N^2}{2\sigma^2}\right)}{2 \pi}\int_{-\infty}^{\infty}\exp\left(-\frac{\left(\sigma t -iy_N/\sigma \right)^2}{2}\right)dt
	\end{align*}


	The integrand is a holomorphic function with no poles\sidenotemark thus we can perform the substitution $z=\sigma t -iy_N/\sigma$:
	
	\begin{align*}
		g(y_N)&=\frac{\exp\left(-\frac{y_N^2}{2\sigma^2}\right)}{2 \pi\sigma}\int_{-\infty}^{\infty}\exp\left(-\frac{z^2}{2}\right)dz\\
		&=\frac{1}{\sqrt{2\pi}\sigma}\exp\left(-\frac{y_N^2}{2\sigma^2}\right) \equiv \N(y_N; 0, \sigma^2)
	\end{align*}\qed

\end{example}
\sidenotetext{If complex analysis isn't your strong suit, don't worry, you can kind of just follow along}

This is the Gaussian (or normal) pdf, and the way we derived it highlights its fundamental importance in statistics. It is important to emphasize that the shape of this limit distribution does not depend on the original distribution of $x$.
One can easily verify that the expected value and variance of the Gaussian coincide with those computed for $y_N$ at the beginning of the demonstration.

In passing, we note that the sum of Gaussian distributions it is still a Gaussian. We will not go into detail about the properties of Gaussian distributions now, since we will encounter this important distribution many times throughout the rest of these notes.

\vspace*{1cm}
\begin{center}
	\textbf{Exercises at: \ref{exer:probability_101}}
\end{center}